\documentclass[11pt]{article}\usepackage[margin=1in]{geometry}\usepackage{amsmath,amsthm}\newtheorem{theorem}{Theorem}\newtheorem{remark}{Remark}
\title{Residual Multichannel Budgets after Dispersion Transfer}\author{}\date{21 July 2026}\begin{document}\maketitle
\begin{abstract}We recompute the normalized defect budget after isolating the
restricted continuous dispersion contribution and auditing the remaining
sector channels.\end{abstract}
Let $q_*$ be the total ceiling and suppose certified channels contribute $q_C<q_*$. For open channels $c$ choose positive shares $\alpha_c$ summing to one.
\begin{theorem}The exact residual allocation is
\[q_c=\alpha_c(q_*-q_C),\qquad F_c\ge \ell_c-\log q_c.
\]
The channel union is at most $q_*$ whenever every certified premise and every open-channel numerator bound holds uniformly over the same retained state.\end{theorem}
\begin{proof}Subadditivity gives the union bound. The residual identity follows from $\sum_c\alpha_c=1$, and substitution into $e^{\ell_c-F_c}\le q_c$ gives the free-energy conditions.\end{proof}
The restricted continuous calculation leaves a positive residual budget for center, cut/orbit, and integer-topology channels.
\begin{remark}This is not yet a physical union bound: the continuous premises and all three remaining channel ratios are still nonuniform or open. The single-star repair is excluded by Stage 44.\end{remark}\end{document}
